\section{Related Literature}
\label{sec:literature}

The problem belongs first to the literature on corruption as a self sustaining order. \citet{andvig_moene_1990}, \citet{bardhan_1997}, \citet{shleifer_vishny_1993}, and \citet{tirole_1996} all study environments in which conduct that is plainly undesirable in welfare terms nevertheless remains privately coherent because each participant takes as given an order sustained by the conduct of others. The version studied here is narrower. The opening arrangement is not an already stratified bureaucracy or a market for fragmented bribes, but a reciprocally corrupt society in which each household both steals and is stolen from, so persistence comes from exact accounting rather than from institutional legitimacy. That emphasis also brings the paper close to the literature on conventions and equilibrium selection. In \citet{young_1993}, local perturbations shift continuation incentives and move societies across self sustaining conventions. The first honest household plays that role here. It does not purify the environment in one step. It changes the local payoff structure of adjacent households, and the resulting windfall and shortfall alter who still finds participation privately worthwhile.

Once that reciprocal order breaks, the relevant comparison is no longer between corruption and honesty in the abstract, but between decentralized predation, hierarchical predation, and organized defense. \citet{olsen_torsvik_1998} are useful because hierarchy there changes both the distribution of rents and the feasibility of corruption itself; the same issue arises once rich households stop stealing personally and begin to purchase predatory labor from poor households. The difference is that hierarchy is not present at the start and has to be generated from the breakdown of a decentralized arrangement. Patronage therefore cannot rest on a vague appeal to elite power. It requires a surplus condition. Delegation has to be profitable because a patron reaches or disciplines targets more effectively than a freelancer can, and because the reservation utility of poor workers remains compressed. That argument then joins the appropriation and defense literature in \citet{skaperdas_1992}, \citet{grossman_kim_1995}, and \citet{hirshleifer_1995}, where insecure claims induce agents to allocate resources between conflict and protection, and it remains close to \citet{demsetz_1967} in treating protection as a costly response to rising stakes. The point of departure matters throughout. The starting environment here is already an order, but one with universal and roughly egalitarian insecurity of claims, so guarding is not the first appearance of order. One order replaces another, and the later order protects a distribution created during the transition itself.

That last step places the paper in the literature on endogenous enforcement and state formation. \citet{olson_1993} distinguishes roving from stationary bandits and argues that future extraction can create incentives to provide order. \citet{north_wallis_weingast_2009} generalize that intuition into violence orders organized around elite management of rents. \citet{besley_persson_2009} study the emergence of legal and fiscal capacity, while \citet{acemoglu_ticchi_vindigni_2011} and \citet{acemoglu_robinson_2008} show how coercive institutions can persist because they serve those who control them. Guards and police belong in that lineage here, because concentrated wealth becomes worth defending and labor can be reallocated from predation to protection without changing who controls the terms of social order. \citet{becker_1968} and \citet{becker_stigler_1974} remain relevant for the same reason. Once guarding becomes available, elite households compare two costly technologies, one extractive and one protective, even though no benevolent state solves an optimal policing problem in the background. The graph theoretic analogy to \citet{shapley_scarf_1974} is useful only at the opening benchmark, where the balanced theft cycle is also a directed cycle with one incoming and one outgoing link per household. It cannot support any stronger claim about core stability or voluntary exchange, because transfers here are coercive, exposure alters later incentives, and rivalry in extraction imposes externalities on agents who have consented to nothing. Taken together, these literatures leave a clear space for the present paper: reciprocal corruption, hierarchical predation, and selective protection appear here as successive forms of one order, and welfare cannot be read from aggregate wealth once the benchmark preserves it by construction. The relevant comparison runs instead through concavity in wealth together with the costs of theft, punishment, contracting, and guarding.
